\documentclass{amsart}
% For dealing with references we use the comment environment
\usepackage{verbatim}
\newenvironment{reference}{\comment}{\endcomment}
%\newenvironment{reference}{}{}
\newenvironment{slogan}{\comment}{\endcomment}
\newenvironment{history}{\comment}{\endcomment}

% For commutative diagrams we use Xy-pic
\usepackage[all]{xy}

% We use 2cell for 2-commutative diagrams.
\xyoption{2cell}
\UseAllTwocells

% We use multicol for the list of chapters between chapters
\usepackage{multicol}

% This is generally recommended for better output
\usepackage{lmodern}
\usepackage[T1]{fontenc}

% For cross-file-references

% Package for hypertext links:
\usepackage{hyperref}

% For any local file, say "hello.tex" you want to link to please

% Theorem environments.
%
\theoremstyle{plain}
\newtheorem{theorem}[subsection]{Theorem}
\newtheorem{proposition}[subsection]{Proposition}
\newtheorem{lemma}[subsection]{Lemma}

\theoremstyle{definition}
\newtheorem{definition}[subsection]{Definition}
\newtheorem{example}[subsection]{Example}
\newtheorem{exercise}[subsection]{Exercise}
\newtheorem{situation}[subsection]{Situation}

\theoremstyle{remark}
\newtheorem{remark}[subsection]{Remark}
\newtheorem{remarks}[subsection]{Remarks}

\numberwithin{equation}{subsection}

% Macros
%
\def\lim{\mathop{\mathrm{lim}}\nolimits}
\def\colim{\mathop{\mathrm{colim}}\nolimits}
\def\Spec{\mathop{\mathrm{Spec}}}
\def\Hom{\mathop{\mathrm{Hom}}\nolimits}
\def\Ext{\mathop{\mathrm{Ext}}\nolimits}
\def\SheafHom{\mathop{\mathcal{H}\!\mathit{om}}\nolimits}
\def\SheafExt{\mathop{\mathcal{E}\!\mathit{xt}}\nolimits}
\def\Sch{\mathit{Sch}}
\def\Mor{\mathop{\mathrm{Mor}}\nolimits}
\def\Ob{\mathop{\mathrm{Ob}}\nolimits}
\def\Sh{\mathop{\mathit{Sh}}\nolimits}
\def\NL{\mathop{N\!L}\nolimits}
\def\CH{\mathop{\mathrm{CH}}\nolimits}
\def\proetale{{pro\text{-}\acute{e}tale}}
\def\etale{{\acute{e}tale}}
\def\QCoh{\mathit{QCoh}}
\def\Ker{\mathop{\mathrm{Ker}}}
\def\Im{\mathop{\mathrm{Im}}}
\def\Coker{\mathop{\mathrm{Coker}}}
\def\Coim{\mathop{\mathrm{Coim}}}

% Boxtimes
%
\DeclareMathSymbol{\boxtimes}{\mathbin}{AMSa}{"02}

%
% Macros for moduli stacks/spaces
%
\def\QCohstack{\mathcal{QC}\!\mathit{oh}}
\def\Cohstack{\mathcal{C}\!\mathit{oh}}
\def\Spacesstack{\mathcal{S}\!\mathit{paces}}
\def\Quotfunctor{\mathrm{Quot}}
\def\Hilbfunctor{\mathrm{Hilb}}
\def\Curvesstack{\mathcal{C}\!\mathit{urves}}
\def\Polarizedstack{\mathcal{P}\!\mathit{olarized}}
\def\Complexesstack{\mathcal{C}\!\mathit{omplexes}}
% \Pic is the operator that assigns to X its picard group, usage \Pic(X)
% \Picardstack_{X/B} denotes the Picard stack of X over B
% \Picardfunctor_{X/B} denotes the Picard functor of X over B
\def\Pic{\mathop{\mathrm{Pic}}\nolimits}
\def\Picardstack{\mathcal{P}\!\mathit{ic}}
\def\Picardfunctor{\mathrm{Pic}}
\def\Deformationcategory{\mathcal{D}\!\mathit{ef}}

\setlength{\emergencystretch}{3em}
\begin{document}
\title[Stacks Project, Tag 091I]{408 --- Strictification of inverse systems of derived objects}
\maketitle
\noindent Copyright (C) 2005--2025 Johan de Jong. Stacks Project authors. Source: \href{https://stacks.math.columbia.edu/tag/091I}{Tag 091I}.

\noindent Editorial difficulty note (not author text). Difficulty H2: The author replaces each derived transition map by a roof and constructs quasi-isomorphic chain-level representatives inductively. The new transition maps must be actual maps of complexes compatible with the existing system, producing a single derived object in the category of module systems..

\noindent Editorial provenance note (not author text). History: excerpt from revision \texttt{a0444\allowbreak 6e57e\allowbreak c1fbc\allowbreak 252a8\allowbreak 71afc\allowbreak ec775\allowbreak 2fb28\allowbreak 07b14}, more-algebra.tex, lines 23681--23750. Reference presentation was adapted; original-author context and core support blocks were retained or added. The mathematical wording is unchanged. Licensed under GNU FDL 1.2 or later; no invariant sections or cover texts. See ../licenses/COPYING. Context blocks reproduce author definitions and supporting results; they are not additional counted problems.

\begin{lemma}
\label{lemma-lift-to-system-complexes}
Let $(A_n)$ be an inverse system of rings. Suppose that we are given
\begin{enumerate}
\item for every $n$ an object $K_n$ of $D(A_n)$, and
\item for every $n$ a map $\varphi_n : K_{n + 1} \to K_n$ of
$D(A_{n + 1})$ where we think of $K_n$ as an object of $D(A_{n + 1})$
by restriction via $A_{n + 1} \to A_n$.
\end{enumerate}
There exists an object
$M = (M_n^\bullet) \in D(\textit{Mod}(\mathbf{N}, (A_n)))$
and isomorphisms $\psi_n : M_n^\bullet \to K_n$ in $D(A_n)$
such that the diagrams
$$
\xymatrix{
M_{n + 1}^\bullet \ar[d]_{\psi_{n + 1}} \ar[r] & M_n^\bullet \ar[d]^{\psi_n} \\
K_{n + 1} \ar[r]^{\varphi_n} & K_n
}
$$
commute in $D(A_{n + 1})$.
\end{lemma}

\begin{proof}
We write out the proof in detail. For an $A_n$-module $T$ we write
$T_{A_{n + 1}}$ for the same module viewd as an $A_{n + 1}$-module.
Suppose that $K_n^\bullet$ is a complex of $A_n$-modules representing $K_n$.
Then $K_{n, A_{n + 1}}^\bullet$ is the same complex, but viewed as a
complex of $A_{n + 1}$-modules. By the construction of the derived
category, the map $\psi_n$ can be given as
$$
\psi_n = \tau_n \circ \sigma_n^{-1}
$$
where $\sigma_n : L_{n + 1}^\bullet \to K_{n + 1}^\bullet$
is a quasi-isomorphism of complexes of $A_{n + 1}$-modules
and $\tau_n : L_{n + 1}^\bullet \to K_{n, A_{n + 1}}^\bullet$
is a map of complexes of $A_{n + 1}$-modules.

\medskip\noindent
Now we construct the complexes $M_n^\bullet$ by induction. As base case
we let $M_1^\bullet = K_1^\bullet$. Suppose we have already constructed
$M_e^\bullet \to M_{e - 1}^\bullet \to \ldots \to M_1^\bullet$
and maps of complexes $\psi_i : M_i^\bullet \to K_i^\bullet$ such that
the diagrams
$$
\xymatrix{
M_{n + 1}^\bullet \ar[d]_{\psi_{n + 1}} \ar[rr] & &
M_{n, A_{n + 1}}^\bullet \ar[d]^{\psi_{n, A_{n + 1}}} \\
K_{n + 1}^\bullet & L_{n + 1}^\bullet \ar[l]_{\sigma_n} \ar[r]^{\tau_n} &
K_{n, A_{n + 1}}^\bullet
}
$$
above commute in $D(A_{n + 1})$ for all $n < e$. Then we consider the diagram
$$
\xymatrix{
& & M_{e, A_{e + 1}}^\bullet \ar[d]^{\psi_{e, A_{e + 1}}} \\
K_{e + 1}^\bullet & L_{e + 1}^\bullet \ar[r]^{\tau_e} \ar[l]_{\sigma_e} &
K_{e, A_{e + 1}}^\bullet
}
$$
in $D(A_{e + 1})$. Because $\psi_e$ is a quasi-isomorphism,
we see that $\psi_{e, A_{e + 1}}$ is a quasi-isomorphism too.
By the definition of morphisms in $D(A_{e + 1})$ we can
find a quasi-isomorphism
$\psi_{e + 1} : M_{e + 1}^\bullet \to K_{e + 1}^\bullet$
of complexes of $A_{e + 1}$-modules such that there exists a morphism
of complexes $M_{e + 1}^\bullet \to M_{e, A_{e + 1}}^\bullet$
of $A_{e + 1}$-modules representing the composition 
$\psi_{e, A_{e + 1}}^{-1} \circ \tau_e \circ \sigma_e^{-1}$
in $D(A_{e + 1})$. Thus the lemma holds by induction.
\end{proof}
\end{document}
